

\section{Frameworks MVC}


\begin{frame}{Le framework Django}

Django est un framework en langage Python pour développer
des applications Web.

\begin{itemize}
   \item S'appuie sur le modèle MVC.
    \item Propose tous les outils concevables: traduction, services, authentification,
        etc, etc.
    \item Repose beaucoup sur les \alert{conventions} de codage (nom des fichiers,
       variables) et le moins possible sur la \alert{configuration}.   
    \item Comprend un moteur de template, un \textit{mapping} vers les bases
        de données, des fonctions de service Web.  
\end{itemize}  

Représentatif d'un framework MVC. Très utilisé pour de bonnes raisons : simple,
puissant, concis, performant.
\end{frame}


\begin{frame}[fragile]{Pour installer}

Il faut disposer d'un environnement Python: \url{https://www.python.org/downloads}{https://www.python.org/downloads}

\vfill

Pour installer Django: il suffit en principe d'exécuter:

\begin{minted}{bash}
pip3 install django
\end{minted}

\vfill

Vérifier que Django est bien installé.

\begin{verbatim}
bash-3.2$ python3
>>> import django
>>> print(django.__path__)
>>> print(django.get_version())
\end{verbatim}

\vfill

À partir de maintenant on va suivre le 
\url{Tutoriel}{https://docs.djangoproject.com/en/1.8/intro/tutorial01/}

\end{frame}


\begin{frame}[fragile]{Créer un nouveau projet}

Django est installé avec un utilitaire \texttt{django-admin}.

\begin{minted}{bash}
django-admin startproject messagerie
\end{minted}

\vfill

On peut lancer un serveur

\begin{minted}{bash}
python3 manage.py runserver
\end{minted}

\vfill

Le serveur est en écoute sur le port 8000: \url{Tester}{http://localhost:8000}

\vfill

À partir de là, on peut travailler sur le projet avec un éditeur
de texte ou (mieux) un IDE comme Eclipse.

\end{frame}

\begin{frame}{Organisation d'un projet}

Un projet Django = un ensemble \alert{d'applications}.

\vfill

\includegraphics[width=12cm]{inidjango}

\vfill

Au départ, une application nommée comme le projet\,: elle
contient la configuration.
\end{frame}

\begin{frame}[fragile]{Créons notre première application}

Avec la commande suivante:
\begin{minted}{bash}
python3 manage.py startapp polls
\end{minted}

\vfill

\includegraphics[width=11cm]{django-newapp}

\vfill

Ensuite, il faut ajouter l'application \texttt{polls} dans \texttt{mysite/settings.py}.
\end{frame}

\begin{frame}[fragile]{Nos modèles}

On définit des \alert{classes} représentant les objets de la base.

\begin{minted}{python}
from django.db import models

class Question(models.Model):
    question_text = models.CharField(max_length=200)
    pub_date = models.DateTimeField('date published')

    def __str__(self):           
        return self.question_text

class Choice(models.Model):
    question = models.ForeignKey(Question)
    choice_text = models.CharField(max_length=200)
    votes = models.IntegerField(default=0)
    
    def __str__(self):           
        return self.choice_text
    
\end{minted}

Django se charge de créer et de maintenir la base de données. 
\end{frame}


\begin{frame}[fragile]{Gestion de la base de données}

Génération des requêtes SQL:

\begin{minted}{bash}
python3 manage.py makemigrations polls
\end{minted}

\vfill

Contrôle des requêtes SQL
\begin{minted}{bash}
python3  manage.py sqlmigrate polls 0001
\end{minted}

\vfill

Exécution des requêtes SQL
\begin{minted}{bash}
python3 manage.py migrate
\end{minted}

\vfill

Django mémorise toutes les évolutions de la base et produit les commandes pour
les exécuter.

\end{frame}


\begin{frame}[fragile]{L'interface d'administration}

\alert{Django a créé une interface Web d'administration des données!}

\vfill
Commençons par créer un super-utilisateur.

\begin{minted}{bash}
python3 manage.py createsuperuser
\end{minted}

\vfill

Indiquer qu'on veut engendrer une interface sur les questions: on ajoute
les lignes suivantes dans \texttt{polls/admin.py}.

\begin{minted}{python}
from django.contrib import admin

from .models import Question

admin.site.register(Question)
\end{minted}

\end{frame}

\begin{frame}{L'interface d'administration}

Automatiquement engendrée et gérée par le framework.
\vfill

\includegraphics[width=11cm]{admin-django}

\vfill

Tester les insertions et mises à jour.

\vfill

Ajouter l'interface sur les Choix, de la même manière.

\end{frame}

\subsection{Les vues}

\begin{frame}[fragile]{Passons aux vues}

La notion de \alert{vue} dans Django correspond plutôt
à la notion de \alert{contrôleur} dans d'autres frameworks.

\vfill

Une vue est constituée \alert{d'actions}. Dans \texttt{polls/views.py}

\begin{minted}{python}
from django.http import HttpResponse

def index(request):
    return HttpResponse("Je suis l'application 'polls', action 'index'.")
\end{minted}

\vfill

Une action reçoit une \alert{requête} (HTTP) et retourne une \alert{réponse} (HTTP).
\end{frame}

\begin{frame}[fragile]{Les actions et leurs URLs}


\alert{Chaque action est associée à une URL}. 

\vfill

Dans \texttt{mysite/urls.py}.

\begin{minted}{python}
from django.conf.urls import include, url
from django.contrib import admin

urlpatterns = [
    url(r'^admin/', include(admin.site.urls)),
    url(r'^polls/', include('polls.urls')),
]
\end{minted}

\vfill

Et dans \texttt{polls/urls.py}.

\begin{minted}{python}
from . import views

urlpatterns = [
    url(r'^$', views.index, name='index'),
]
\end{minted}

\end{frame}


\begin{frame}[fragile]{Dernière étape: les templates}

Template = un fragment HTML, avec des commandes pour afficher un 
\alert{contexte.}

\vfill

Créer le fichier \texttt{polls/templates/polls/index.html}.

\begin{minted}{html}
{% if questions %}
    <ul>
    {% for question in questions %}
        <li>
             {{ question.question_text }}
        </li>
    {% endfor %}
    </ul>
{% else %}
    <p>No polls are available.</p>
{% endif %}
\end{minted}

\vfill

"Contexte" = des données passées au template par l'action.
\end{frame}

\begin{frame}[fragile]{L'action crée le contexte et appelle le template}

Redéfinir l'action \texttt{index} dans \texttt{polls/views.py} de la manière suivante:

\begin{minted}{python}
from .models import Question

def index(request):
    questions = Question.objects.order_by('-pub_date')[:5]
    context = {'questions': questions}
    
    return render(request, 'polls/index.html', context)
\end{minted}

\vfill

tester!

\end{frame}

\begin{frame}{Récapitulons}

\textit{Framework} = cadre de travail. C'est bien ça, on développe une application
en définissant

\begin{itemize}
  \item Des \alert{modèles}, qui représentent les données et la logique "métier".
  \item Des \alert{actions} qui traitent des requêtes HTTP.
  \item Des \alert{templates} (ou vues) qui affichent.
\end{itemize}

\vfill
Et tout un environnement complémentaire: tests, déploiement, intégration, etc.


\vfill

Pour aller plus loin: reprendre le tutoriel Django (qui est plus complet) et le faire
dans son intégralité.



\end{frame}